% GENERATED FILE. Edit canonical lessons, then run npm run generate:books.
% canonical-source-hash: fnv1a64:8248d669c20623e4
% canonical-lessons: KA-C131-havu, KA-C131-ame, KA-C131-gili, KA-C131-parivala, KA-C131-gube

\chapter{Snake, Tortoise, Parrot, Pigeon, Owl}
\label{ch:ka-snake-tortoise-parrot-pigeon-owl}
\hlchaptermodality{131}

\begin{chapteropening}
\textbf{By the end of this chapter:} \emph{I can say a snake, a tortoise, a parrot, a pigeon and an owl.}

Complete the last lesson of chapter 131: I can say a snake, a tortoise, a parrot, a pigeon and an owl.
\end{chapteropening}

\section[\texorpdfstring{\kn{ಹಾವು} (hāvu)}{ಹಾವು (hāvu)}]{\kn{ಹಾವು} (hāvu) — a snake}

\label{lesson:KA-C131-havu}



\begin{quote}
Before the new one: say the Kannada for a deer, then the Kannada for a fox.
\end{quote}

\subsection*{You'll want to know: \kn{ಹಾವು}}
\textbf{\kn{ಹಾವು}} — \emph{hāvu} — \textquotedblleft{}a snake\textquotedblright{}.

\begin{grammarlens}[title={What you've built}]
One more word for this chapter.
\end{grammarlens}

\subsection*{Your turn}
\begin{itemize}
\raggedright
  \item \emph{Say it:} \emph{hāvu}
  \item \emph{Say it:} \emph{hāvu}, once more
  \item \emph{Say it:} say \emph{hāvu}
  \item \emph{Recall:} say the Kannada for a deer, then the Kannada for a fox, then say \emph{hāvu} again
\end{itemize}

\begin{tcolorbox}[breakable,colback=teal!4,colframe=teal!35!black,title={Before you move on}]
What does \kn{ಹಾವು} mean? (\textquotedblleft{}A snake\textquotedblright{}.) Say it once more.
\end{tcolorbox}
\section[\texorpdfstring{\kn{ಆಮೆ} (āme)}{ಆಮೆ (āme)}]{\kn{ಆಮೆ} (āme) — a tortoise}

\label{lesson:KA-C131-ame}



\begin{quote}
Before the new one: say the Kannada for a fox, then the Kannada for a snake.
\end{quote}

\subsection*{You'll want to know: \kn{ಆಮೆ}}
\textbf{\kn{ಆಮೆ}} — \emph{āme} — \textquotedblleft{}a tortoise\textquotedblright{}.

\begin{grammarlens}[title={What you've built}]
One more word for this chapter.
\end{grammarlens}

\subsection*{Your turn}
\begin{itemize}
\raggedright
  \item \emph{Say it:} \emph{āme}
  \item \emph{Say it:} \emph{āme}, once more
  \item \emph{Say it:} say \emph{āme}
  \item \emph{Recall:} say the Kannada for a fox, then the Kannada for a snake, then say \emph{āme} again
\end{itemize}

\begin{tcolorbox}[breakable,colback=teal!4,colframe=teal!35!black,title={Before you move on}]
What does \kn{ಆಮೆ} mean? (\textquotedblleft{}A tortoise\textquotedblright{}.) Say it once more.
\end{tcolorbox}
\section[\texorpdfstring{\kn{ಗಿಳಿ} (giḷi)}{ಗಿಳಿ (giḷi)}]{\kn{ಗಿಳಿ} (giḷi) — a parrot}

\label{lesson:KA-C131-gili}



\begin{quote}
Before the new one: say the Kannada for a snake, then the Kannada for a tortoise.
\end{quote}

\subsection*{You'll want to know: \kn{ಗಿಳಿ}}
\textbf{\kn{ಗಿಳಿ}} — \emph{giḷi} — \textquotedblleft{}a parrot\textquotedblright{}.

\begin{grammarlens}[title={What you've built}]
One more word for this chapter.
\end{grammarlens}

\subsection*{Your turn}
\begin{itemize}
\raggedright
  \item \emph{Say it:} \emph{giḷi}
  \item \emph{Say it:} \emph{giḷi}, once more
  \item \emph{Say it:} say \emph{giḷi}
  \item \emph{Recall:} say the Kannada for a snake, then the Kannada for a tortoise, then say \emph{giḷi} again
\end{itemize}

\begin{tcolorbox}[breakable,colback=teal!4,colframe=teal!35!black,title={Before you move on}]
What does \kn{ಗಿಳಿ} mean? (\textquotedblleft{}A parrot\textquotedblright{}.) Say it once more.
\end{tcolorbox}
\section[\texorpdfstring{\kn{ಪಾರಿವಾಳ} (pārivāḷa)}{ಪಾರಿವಾಳ (pārivāḷa)}]{\kn{ಪಾರಿವಾಳ} (pārivāḷa) — a pigeon}

\label{lesson:KA-C131-parivala}



\begin{quote}
Before the new one: say the Kannada for a tortoise, then the Kannada for a parrot.
\end{quote}

\subsection*{You'll want to know: \kn{ಪಾರಿವಾಳ}}
\textbf{\kn{ಪಾರಿವಾಳ}} — \emph{pārivāḷa} — \textquotedblleft{}a pigeon\textquotedblright{}.

\begin{grammarlens}[title={What you've built}]
One more word for this chapter.
\end{grammarlens}

\subsection*{Your turn}
\begin{itemize}
\raggedright
  \item \emph{Say it:} \emph{pārivāḷa}
  \item \emph{Say it:} \emph{pārivāḷa}, once more
  \item \emph{Say it:} say \emph{pārivāḷa}
  \item \emph{Recall:} say the Kannada for a tortoise, then the Kannada for a parrot, then say \emph{pārivāḷa} again
\end{itemize}

\begin{tcolorbox}[breakable,colback=teal!4,colframe=teal!35!black,title={Before you move on}]
What does \kn{ಪಾರಿವಾಳ} mean? (\textquotedblleft{}A pigeon\textquotedblright{}.) Say it once more.
\end{tcolorbox}
\section[\texorpdfstring{\kn{ಗೂಬೆ} (gūbe)}{ಗೂಬೆ (gūbe)}]{\kn{ಗೂಬೆ} (gūbe) — an owl}

\label{lesson:KA-C131-gube}



\begin{quote}
Before the new one: say the Kannada for a parrot, then the Kannada for a pigeon.
\end{quote}

\subsection*{You'll want to know: \kn{ಗೂಬೆ}}
\textbf{\kn{ಗೂಬೆ}} — \emph{gūbe} — \textquotedblleft{}an owl\textquotedblright{}.

\begin{grammarlens}[title={What you've built}]
One more word for this chapter.
\end{grammarlens}

\subsection*{Your turn}
\begin{itemize}
\raggedright
  \item \emph{Say it:} \emph{gūbe}
  \item \emph{Say it:} \emph{gūbe}, once more
  \item \emph{Say it:} say \emph{gūbe}
  \item \emph{Recall:} say the Kannada for a parrot, then the Kannada for a pigeon, then say \emph{gūbe} again
\end{itemize}

\begin{tcolorbox}[breakable,colback=teal!4,colframe=teal!35!black,title={Before you move on}]
What does \kn{ಗೂಬೆ} mean? (\textquotedblleft{}An owl\textquotedblright{}.) Say it once more.
\end{tcolorbox}
